
\section{Dynamic Programming for Odd Powers}
\label{sec:dp}

The approach to odd powers in Section~\ref{sec:odd} involves fully covering cosets using Methods 1 and 2. In this section, we permit partial coverage, allowing the available positions and symbol sets to be distributed more efficiently between cosets. This results in a computable lower bound for $M(q+1,q)$, and in certain cases this is an improvement over the result in Section~\ref{sec:odd}.

\subsection{Partial Coverage of a Single Coset}


We give a dynamic programming (DP) algorithm that yields lower bounds that are better than those given by Theorem \ref{odd}. In the technique used for that theorem, recall that we fully covered cosets by using either:

{\bf Method 1}: take the set of values to be an entire suffix set, say $Suffix(\alpha)$, and the set of positions to be one stripe (i.e., one permutation from each suffix set).

{\bf Method 2}: take the set of values to be one of the special sets and the set of positions be $p$ stripes (i.e., one choice from each special set).

Now we will consider a generalized method which may cover less than the whole coset:

{\bf General Method}: take the set of values to be the union of $b$ of the special sets and the set of positions be $a$ stripes (i.e., one choice from each suffix set for $a$ different prefixes).

You see that if $a=1$ and $b=p$, then the General Method is equivalent to Method 1. Similarly, if $a=p$ and $b=1$, then the General Method is equivalent to Method 2. However, when $a$ and $b$ are different values, then this method may cover less than the whole coset.

\begin{theorem}
\label{th:general coset coverage}

Using the General Method, a coset's coverage is:
$$ p^{2m} * |\{y-x \pmod p : x \in P', y \in Q'  \}|  $$
Where $P'$ is the set of prefix sums represented by the $a$ stripes in $P$, and $Q'$ is the set of prefix sums represented by the $b$ special sets in $Q$.

\end{theorem}

\begin{proof}

(TODO, change this to use shift())
Since $P'$ and $Q'$ are sets of prefixes, let's define the set of prefixes $V = \{ y-x | x \in P', y \in Q' \}$. For some coset $i$, we cover permutation $j$ when $y = i*x+j$ has a solution for $x \in P_i$ and $y \in Q_i$. In other words, we cover any $j = y - i*x$. In other words, we cover permutation $j$ when $prefix(j) \in V_i$. So in coset $i$, we cover exactly $|V_i| * p^{2m}$ permutations.

\end{proof}

\subsection{Allocating Positions and Symbols Between Multiple Cosets}

The General Method described above uses a set of $a*p^m$ positions and a set of $b*p^m$ symbols. There are a total of $p^{2m+1}$ elements in $GF(p^{2m+1})$, which are considered as both positions and symbols. Therefore, if we want to maximize our coverage of multiple cosets using the General Method, we must find a balance between getting enough coverage on each coset while saving enough positions and symbols to use on the other cosets.

This optimization is not straightforward. To begin, we make some simplifying assumptions. First, instead of trying to construct the final permutation array, we are only going to construct a property $I$, namely that for some coset $i$, $P_i$ contains $a$ stripes and $Q_i$ contains $b$ special sets. Further, that the $a$ stripes of $P_i$ have prefix sums that are in consecutive order, and the $b$ special sets have prefix sums that are in consecutive order skipping $t$ at a time. Lastly, that when $Q_i$ has some skipping factor $t$, then the same $t$ must be used for $t$ cosets in a row. In other words, if $Q_0$ has $m=3$, then $Q_1$ and $Q_2$ must also have $t=3$.

From this, we have three sequences of numbers to optimize:
- $A_s$ is the number of stripes given to each set of positions.
- $B_s$ is the number of $special~sets$ given to each set of symbols.
- $T_s$ is the skipping factor, and hence also number of cosets being assigned in a row

We also have two budgets to account for. The first budget is the number of special sets available for each $Q_i$, and this budget starts at a value of $p^{m+1}$. The second budget is the number of stripes available to make each $P_i$, which also starts at $p^{m+1}$ since there are $p$ different prefixes and each special set is of size $p^m$.

Applying Theorem \ref{th:general coset coverage} to these budgets and sequences, we find that we are trying to maximize this:
$$ p^{2m} \sum_{(a,b,t)} |\{yt-x \pmod p : 0 \le x < a,\ 0 \le y < b\}|$$

Dropping the $p^{2k}$ term for brevity, we convert this into a recurrence for dynamic programming:
\begin{itemize}

\item  $\theta (0,0) = 0$

\item $\theta (p,q) = max\left\{ \theta(p-at, q-bt) + |\{yt_s-x \pmod p : 0 \le x < a_s,\ 0 \le y < b_s\}| \right\}$

\end{itemize}

\subsection{Pseudocode and Analysis}

\begin{minted}[
    frame=lines,
    framesep=2mm,
    baselinestretch=1.2,
    fontsize=\footnotesize,
    linenos
]{text}

Algorithm dp_split(p,m)

Input:
  prime p and integer m > 0

  B      ← p^(m+1)
  A      ← (B+1)*(B+1) sized array, initialized to all zeros  
  choice ← (B+1)*(B+1) sized array, initialized to all null  

  for u in 1..B:
    for v in 1..B:
      for a in 1..min(p, u):
        for b in 1..min(p, v):
          limit ← min(p, floor(u/a), floor(v/b))
          for t in 1..limit:
            gain ← t * overlap(p, a, b, t)
            candidate ← gain + A[u-t*a, v-t*b]
            if candidate > A[u,v]:
              A[u,v] ← candidate
              choice[u,v] ← (a, b, t)

  split ← empty list
  u = B
  v = B
  while choice[u,v] not null:
    a, b, t = choice[u,v]
    append (a,b,t) to split
    u ← u-a*t
    v ← v-b*t
  return A[B,B], split


\end{minted}

We see from the size of the variables that the space complexity is $O(p^{2m+2})$. From the nested `for` loops, we see that the complexity of building the table is $O(p^{2m+3})$.

The correctness of this algorithm follows from the optimal substructure. 

\subsection{Constructing a PA from the Dynamic Programming Output}

\noindent
Let
\[
\mathcal S=((a_1,b_1,t_1),\ldots,(a_s,b_s,t_s))
\]
be the allocation sequence returned by Algorithm~\ref{alg:dp}, and let
\[
K_r=\sum_{j<r} t_j.
\]
Construct a bipartite graph $G=(A,B,E)$ as follows:
\[
\begin{aligned}
A
  &= \bigcup_{r=1}^{s}
     \left\{
       (K_r+\ell,g,h):
       0\le \ell<t_r,\;
       0\le g<a_r,\;
       0\le h<p^m
     \right\},\\
B
  &= GF(q).
\end{aligned}
\]

For $g\in\{0,\ldots,p-1\}$ and
$h\in\{0,\ldots,p^m-1\}$, let $\mathcal S_{g,h}$ denote
the special set whose prefix is $g$ and whose suffix is $h$.
For each vertex $v=(k,g,h)\in A$, define
\[
N(v)=
\left\{
shift(k,x):x\in\mathcal S_{g,h}
\right\},
\]
Finally, let
\[
E=\{\{v,x\}:v\in A,\ x\in N(v)\}.
\]
